\section{Social Impact \& Conclusion}
\label{sec:conclusion}


Human communication is multisensorial—we take in sensory input from several modalities to process information in a dynamic way \citep{holler2019multimodal}. In multilingual contexts, advancements in text-based machine translation have given rise to tools that help individuals communicate and learn in languages where proficiency is low \citep{lee2023effectiveness}. That said, while foundational models such as NLLB \citep{nllb2022} push \mt{} beyond 200 languages, direct speech translation has yet to achieve similar strides. To bridge this gap, we created a massively multilingual and multimodal machine translation system that paves the way for the next generation of speech translation technologies.

Using novel data and modeling approaches to combine \sst{}, \ttst{}, \st{}, \mt{}, and \asr{} in a single model, our main contributions are as follows. First, we built a new LID model aligned with our language coverage and conducted speech mining with the help of the newly conceived SONAR—a multilingual and multimodal sentence embedding space—to create a corpus of automatically aligned speech translations of more than \totalminedhours hours. By fusing four building blocks, (1) \nllbv, a massively multilingual T2TT model, (2) w2v-BERT 2.0, a speech
representation learning model pre-trained on unlabeled speech audio data, (3) T2U, a text-to-unit sequence-to-sequence model, and (4) HiFi-GAN—a multilingual vocoder for synthesizing
speech from units, we built a unified model that covers \sst{} from \nsslangs{} languages to English (\nsslangs{}-eng), English to \ntslangs{} languages (eng-\ntslangs{}), and \st{} for \nsslangs{}-eng and eng-\ntextlangs{} languages. Notably, compared to previous work on \sst{}, which primarily serves translations into English and not vice versa, \mfourt is capable of performing translation from English towards \ntslangs{} directions. When it comes to \st{}, \mfourt achieves an improvement of 20\% BLEU over the previous state-of-the-art in \st{} translation. Preliminary human evaluations of \st{} outputs evinced similarly impressive results; for translations from English, XSTS scores for 24 evaluated languages are consistently above 4 (out of 5). For into English directions, we see significant improvement over \whisperlarge's baseline for 7 out of 24 languages.
We then evaluated our model for robustness, revealing that \mfourt is more robust than \citep{whisper} when it comes to background noises and
speaker variations. By also including results of the level of added toxicity and gender bias, we hope to motivate future work targeting mitigation efforts.

Made with the goal of promoting accessibility, we open-source all contributions of our work, including two sizes of our model to ensure that even researchers with limited computing resources can use our work. In the section below, we discuss the potential social impact of \mfourt by focusing on its downstream possibilities.

\subsection{Augmenting world-readiness}

The world we live in has never been more interconnected—the global proliferation of the internet, mobile devices, communicative platforms, and social media exposes individuals to more multilingual content than ever before \citep{zuckerman_2008}. The current social order places a demand on a person’s “world-readiness” \citep{ACTFL}, a measure of how competent a person is to take on the polyglot world. Initially developed in the context of language learning, world-readiness underscores the importance of being able to communicate in languages beyond one’s mother tongue for both instrumental (i.e., employment or schooling) and cultural reasons (i.e., to become a global citizen). That said, while we believe that language acquisition should remain a key mechanism for boosting one’s world-readiness, we acknowledge that doing so requires mental and material resources many people may not possess.

The downstream applications that \mfourt supports could allow on-demand access to world-readiness by streamlining multilingual exchange across various contexts. Akin to what \mt{} has accomplished for bridging the comprehension of multilingual texts, \mfourt could have the same effects for speech. Research shows that contrary to one’s native language, where speech is more organically acquired than reading or writing \citep{liberman1992relation}, this tendency is flipped when it comes to foreign languages \citep{cheng1999language}. In other words, speech is often deemed more challenging than reading or writing in a foreign language context. \mfourt-supported applications could act as a co-piloting mechanism that supports users in multilingual conversations and boost their confidence in speech-heavy interactions. As speech-based interfaces (i.e., audio assistants, voice memos, live transcriptions, etc.) and auditory content (i.e., podcasts, audiobooks, short-form videos, etc.) become ever more present in people’s lives, downstream applications enabled by \mfourt could allow a greater variety of multilingual experiences and in ways that feel more natural and dynamic than its text-based counterparts.

From an inclusion standpoint, \mfourt’s focus on multimodality could make a meaningful difference in augmenting the world-readiness of those with accessibility needs and those whose languages contain multiple writing systems (as aforementioned in \ref{sec:appendix:problem}. For many who lack reading or writing skills, or are unable to rely on sight (i.e., people who are blind or with visual impairment), voice-assisted technologies are essential to how they communicate and stay connected \citep{belekar2020voice}. The ability to translate speech not only gives these groups more comprehensive access to information beyond their native languages, but also in a manner that is better suited for their communicative needs. Additionally, recognizing that some languages may have script variance, \mfourt’s offers up affordances that help circumvent the multiscript conundrum. For languages that do not have standardized writing systems, investments in speech recognition and translation may be instrumental in preventing endangerment. We hope that our effort can help contribute to this important movement.

\subsection{Future work}

As is the case with most technologies, the distribution of benefits varies based on user demographics and social situation \citep{wang2023human}. While we make the case that \mfourt could augment world-readiness by lowering the barriers in cross-lingual communication, some users may experience more difficulties using our work than others. For instance, like many other speech technologies, \mfourt’s \asr{} performance may vary based on gender, race, accent, or language \citep{koenecke2020racial,ngueajio2022hey}. Moreover, our system's performance when it comes to translating slang or proper nouns may also be inconsistent across high and low-resource languages. 

Another challenge for \sst{} is that speech hinges on immediate reception and feedback compared to written language. In other words, a speaker is limited in their ability to ascertain the quality of an output or make “edits” in a live conversation. Without the ability to plan and revise with the help of back-translation or a native speaker, \sst{} may carry higher degrees of interactional risks when it comes to mistranslations or toxicity. We urge researchers and developers who fine-tune or build products using \mfourt to think critically about design features that could help users circumvent these potential obstacles. On a related note, we believe that \mfourt-fueled applications should best be viewed as an augmentation device that assists in translation rather than a tool that replaces the need for language learning or reliable human interpreters. This reminder is especially pertinent in high-stakes situations involving legal or medical decision-making.

Finally, speech is not spoken text—it encompasses a suite of prosodic (i.e., rhythm, stress, and intonation) and emotional components that deserve further research \citep{elbow1985shifting}. To create \sst{} systems that feel  organic and natural, more research should be directed at output generation that preserves expressivity \citep{trilla2012sentence}. In addition, the consummate realization of the Babel Fish requires deeper investments into research on low-latency speech translation. Developing systems that enable streaming (i.e., incrementally translating an input sentence as it is being presented) may increase the adoption of such systems in industry or educational contexts \citep{iranzo2022simultaneous,rybakov2022streaming}. We hope that \mfourt opens up new possibilities for both of these research areas. 


